\UTIResumenPagina
  {english}
  {Abstract}
  {Language model inference consists of two phases with different bottlenecks: prefill, which processes the prompt, and decode, which generates the response one token at a time. This work set out to characterize both phases for Qwen3-14B served with vLLM on a 40~GB NVIDIA A100 GPU of the CEDIA HPC cluster, and to measure the effect of quantization, context length and concurrency. A trial-and-error methodology was followed over eleven iterations, each designed from the failure observed in the previous one. None of the four \texttt{vllm serve} start-ups, one with Qwen3-8B and three with Qwen3-14B, became ready: importing the environment from the shared NFS storage grew from 12.9 to 25.8~minutes across attempts, and copying it to the local disk did not complete. Consequently, no Qwen3-14B metrics were obtained. The phases were characterized with Ollama on the full A100 for two 30B models quantized in Q4\_K\_M. For \texttt{nemotron-3-nano:30b}, the time per output token varied by 0.95~\% between 170 and 2390 prompt tokens, while prefill speed increased by a factor of 1.99; for \texttt{qwen3-coder:30b}, the former rose by 7.1~\% and the latter by 56.6~\%. This asymmetry is consistent with a decode phase bound by memory bandwidth and a prefill phase that amortizes weight reads. It was also found that, in Ollama, the client-side time to first token includes a \texttt{load\_duration} term that reaches 94~\% of its value, so prefill must be taken from the server. The general objective was not achieved for the planned combination; the work contributes a characterization of the asymmetry between phases, evidence that start-up from NFS is not viable on this cluster, and a validity criterion for measuring with Ollama.}
  {Keywords}
  {decode, LLM inference, Ollama, prefill}
